%===============================================================================
% SEÇÃO 1 — Introdução
%===============================================================================

\section{Introdução}

A reforma a seco de biogás (BDR, do inglês \textit{biogas dry reforming})
consiste na conversão de uma mistura CH\(_4\)/CO\(_2\) proveniente de digestão
anaeróbia em gás de síntese (H\(_2\) + CO), utilizando o próprio CO\(_2\) do
biogás como agente oxidante do metano \citep{lavoie2014}. AO CITAR REFERENCIA, O PONTO TEM QUE VIR DEPOIS E NAO ANTES COMO ESTAVA

Ao substituir o metano
puro pelo biogás, o processo passa a atuar sobre uma corrente renovável e de
menor valor agregado, contribuindo tanto para a valorização de resíduos orgânicos
quanto para a mitigação de emissões de gases de efeito estufa.\citep{lavoie2014}
Nesse contexto, a BDR se destaca como rota promissora para a produção de
hidrogênio renovável e gás de síntese integrados a cadeias de energia de baixo
carbono.

Do ponto de vista termodinâmico, a reação é fortemente endotérmica, com entalpia
padrão da ordem de 247~kJ\(·\)mol\(^{-1}\), de forma que temperaturas elevadas e
suprimento intenso de calor são necessários para se atingir conversões
significativas.\citep{thermodynamics2025} Estudos de equilíbrio indicam que,
para misturas representativas de biogás, a conversão de metano aumenta
continuamente com a temperatura, situando-se em faixas típicas de operação
entre 650 e 900~\(^\circ\)C, enquanto abaixo de cerca de 640~\(^\circ\)C a
reação deixa de ser espontânea e a formação de carbono sólido torna-se
termodinamicamente favorecida.\citep{thermodynamics2025,lavoie2014} O uso de
biogás, em comparação ao metano puro, traz complexidades adicionais: a
composição CH\(_4\)/CO\(_2\) varia ao longo do tempo e o gás pode conter
espécies inertes (N\(_2\), O\(_2\)) e traços de impurezas como H\(_2\)S, afetando
tanto o desempenho catalítico quanto o comportamento dinâmico do processo.

Apesar do volume expressivo de publicações sobre desenvolvimento de catalisadores,
cinética reacional e análises termodinâmicas para reforma a seco, há relativamente
poucos trabalhos dedicados ao controle e à automação de processos de BDR em
plantas piloto reais \citep{luyben2014,essien2019}. Na maioria dos casos, as
estratégias de controle são avaliadas em reatores de bancada ou em modelos
simulados de DRM com metano puro, sem considerar as flutuações de composição e
as incertezas inerentes ao biogás. Em processos fortemente endotérmicos e
suscetíveis à formação de coque, como a BDR, o controle passa a ser um elemento
central para viabilizar operação estável e segura, especialmente quando se
busca operar em janelas de alta conversão e baixa deposição de
carbono \citep{luyben2014}.

MUDEI ESTE PARAGRAFO

No Brasil, a Universidade Federal do Paraná (UFPR) vem desenvolvendo a rota de reforma a seco de biogás de forma sistemática, deste a escala de bancada até uma unidade piloto.
Trabalhos iniciais demonstraram a viabilidade catalítica do processo,
incluindo a validação de catalisadores à base de níquel e o escalonamento da
síntese catalítica para as quantidades exigidas pelos ensaios em maior
escala \citep{oliveira2021,oliveira2022}. Nessas unidades, as variáveis de processo
(temperatura, vazão e composição) são manipuladas manualmente ou por
controladores simples de laboratório, o que é viável em vista da pequena escala
e da supervisão direta do pesquisador durante os ensaios \citep{oliveira2021}.

A transição da escala de bancada para uma planta piloto operacional, a planta B2H2

%1.
%OLIVEIRA, L. G. ; OLIVEIRA, G. H. C. ; FERREIRA, J. C. S. ; VILELA JUNIOR, J. A. %; ALVES, H. J. . Projeto B2H2: Produção de hidrogênio renovável para geração %distribuída assistida por IA. In: 4o Congresso Brasileiro do Hidrogênio, 2025, %Brasilia. Anais do 4o Congresso Brasileiro do Hidrogênio, 2025.

%2.
%FERREIRA, J. C. S. ; ESCRIBANO, R. A. N. G. ; Schreiner, M. A. ; Ponciano, P. C. %; OLIVEIRA, G. H. C. ; VILELA JUNIOR, J. A. ; ALVES, H. J. . Do Resíduo Orgânico %à Rede Elétrica: Produção On-Grid de Hidrogênio Renovável com Supervisão %Inteligente na Planta B2H2. In: 4o Congresso Brasileiro do Hidrogênio, 2025, %Brasilia. Anais do 4o Congresso Brasileiro do Hidrogênio, 2025.

Quando falar da planta, cite estes dois artigos acima!

, que é uma etapa necessária antes da implantação industrial, impõe uma mudanças qualitativa nos requisitos de operação. A supervisão manual já não é
viável, as distâncias térmicas são maiores, as inércias são mais lentas e os
desvios operacionais têm consequências críticas do ponto de vista de
segurança e custo. Nessa escala, a automação e o controle
de processos deixam de ser acessórios e passam a ser requisitos fundamentais
para manter a operação dentro da janela segura, garantir alta conversão de
CH\(_4\)/CO\(_2\) e proteger o catalisador de deposição acelerada de coque.
É precisamente nesse contexto que se insere o presente trabalho. A planta
piloto B2H2, com reação A SECO (QUAL) operando entre 650--800~\(^\circ\)C e produzindo ao final cerca
de 30~L\(·\)h\(^{-1}\) de hidrogênio em regime contínuo,
constituí a plataforma experimental sobre a qual se formula, implementa e
valida a arquitetura de controle industrial que viabiliza esse \emph{scale-up}
de forma segura e reprodutível.

ESTE LANCE AQUI DO BRD DEVERIA TER IDO PARA ANTES DE COMECAR A FALAR DA UFPR

A BDR é particularmente suscetível à deposição de carbono (coque) sobre o
catalisador, decorrente de reações como a decomposição do metano e a reação
de Boudouard, resultando em perda de atividade catalítica, aumento da perda de
carga e, nos casos mais severos, riscos operacionais por obstrução do
leito.\citep{lavoie2014,thermodynamics2025} No biogás, a variabilidade da
razão CO\(_2\)/CH\(_4\) e a presença eventual de impurezas tornam ainda mais
necessário manter as variáveis de operação dentro de uma janela que faça
prevalecer a reação desejada sobre os caminhos de formação de
coque.\citep{thermodynamics2025,oliveira2021}

O restante deste artigo está organizado da seguinte forma. A
Seção~\ref{sec:formulacao} apresenta a formulação matemática do problema de
controle. AS SECOES 2 e 3 ME PARECEM SER OU TRATAR DA MESMA COISA. TALVEZ UMA PUDESSE SER ELIMINADA. A Seção~\ref{sec:revisao} revisa conceitos de controle de processos
endotérmicos e o uso de análises termodinâmicas como suporte à definição de
janelas operacionais QUE JANELAS VC SE REFERE AQUI?. A Seção~\ref{sec:processo} descreve a configuração da
planta piloto. A Seção~\ref{sec:validacao} apresenta e discute os resultados
experimentais. Por fim, a Seção~\ref{sec:conclusao} traz as conclusões e
indica trabalhos futuros.
